Man vergleicht die Abstände $\Delta s$ aufeinanderfolgender Bilder. Bei einer gleichförmigen Bewegung sind sie gleich. Bei einer gleichmässig beschleunigten Bewegung nimmt die Geschwindigkeit in jedem Intervall um gleich viel zu, die Abstände wachsen also jedes Mal um denselben Betrag $\Delta(\Delta s) = a\,\tau^2$.
\begin{center}
\begin{tabular}{l|cccc|ccc}
 & \multicolumn{4}{c|}{$\Delta s$ in \si{\meter}} & \multicolumn{3}{c}{$\Delta(\Delta s)$ in \si{\meter}} \\
\hline
A & $2$ & $2$ & $2$ & $2$ & $0$ & $0$ & $0$ \\
B & $1$ & $2$ & $3$ & $4$ & $1$ & $1$ & $1$ \\
C & $1$ & $2$ & $4$ & $8$ & $1$ & $2$ & $4$
\end{tabular}
\end{center}
\begin{itemize}
\item A ist $\result{\text{gleichförmig}}$ mit
\begin{align}
 v &= \vAF = \frac{\dsA}{\dT} = \vA \approx \result{\vAP{0}{1}}
\end{align}
\item B ist $\result{\text{gleichmässig beschleunigt}}$: Die Abstände wachsen jedes Mal um \dsBaO. Mit $\Delta(\Delta s) = a\,\tau^2$:
\begin{align}
 a &= \aBF = \frac{\dsBb-\dsBa}{\left(\dT\right)^2} = \aB \approx \result{\aBP{0}{1}}
\end{align}
\item C ist $\result{\text{weder noch}}$: Der Körper wird immer schneller, aber die Abstände wachsen jedes Mal um mehr (sie verdoppeln sich). Die Geschwindigkeit nimmt immer schneller zu; die Beschleunigung ist nicht konstant, sondern wird grösser.
\end{itemize}
